\section{Proofs for Section~\ref{sec:playbook}}\label{app:proofs-playbook}

Notation is as in Section~\ref{sec:playbook} and Part~I. Write $y:=X_{\mathsf o}-k>0$, $\beta=1/\alpha$ and $e:=1/(1-\alpha)$. From \eqref{eq:Dstar}--\eqref{eq:Rstar},
\begin{equation}\label{eq:pb-closed}
\sigma^*=(\alpha\ell y)^{e},\qquad \Delta^*=\ell^{e}(\alpha y)^{\alpha e},\qquad \Pi=(1-\alpha)y\Delta^*=C_\alpha(\ell y)^{e},\qquad \sigma^*=\alpha y\Delta^* ,
\end{equation}
so $y\Delta^*=\sigma^*/\alpha$ and $\Pi=(1-\alpha)\sigma^*/\alpha$.

\paragraph{Proof of Lemma~\ref{lem:playbook-calib}.}
(a) With $\alpha(\hat X-k)=\mathsf r_\sigma$, $\Delta=\sigma/\mathsf r_\sigma$ and $\ell=\Delta\sigma^{-\alpha}$, \eqref{eq:pb-closed} gives $\sigma^*(\ell,\hat X)=(\ell\,\mathsf r_\sigma)^{e}=(\Delta\sigma^{-\alpha}\sigma/\Delta)^{e}=(\sigma^{1-\alpha})^{e}=\sigma$, and $\Delta^*=\sigma^*/(\alpha(\hat X-k))=\sigma/\mathsf r_\sigma=\Delta$; for $\sigma=0$, $\ell=0$ gives $\sigma^*=0=\sigma$ and $\Delta^*=0=\Delta$. (b) By \eqref{eq:pb-closed}, $\sigma^*$ is proportional to $y^{e}$ and $\Delta^*$ to $y^{\alpha e}$ at fixed $\ell$; divide by their values at $\hat y=\hat X-k$ and use (a). (c) $\hat X-k=\mathsf r_\sigma/\alpha$ is decreasing in $\alpha$. With the netted $\mathsf r_\sigma=\SpendRatioPB$ the interval is $[\pbXhatKNetLo,\pbXhatKNetHi]$, and with the all-S\&M ratio of \SpendRatioPBAll{} it is $[\pbXhatKAllLo,\pbXhatKAllHi]$. \qed

\paragraph{Proof of Proposition~\ref{prop:playbook-rank}.}
(a) \emph{Growth engine.} By Proposition~\ref{prop:envelope}, $\partial R^*_X/\partial\lambda=y\Delta^*/\lambda$, so a proportional change $d\ln\lambda$ raises $R^*_X$ by $y\Delta^*\,d\ln\lambda=(\sigma^*/\alpha)\,d\ln\lambda$. A gain $x_\lambda$ in demand pull has $d\ln\lambda=x_\lambda$, which gives $\Delta R_\lambda=x\sigma^*/\alpha$ (Corollary~\ref{cor:priorities}(ii)). A saving $\delta_\sigma=x$ on growth activities has $d\ln\lambda=-\alpha\ln(1-x)=\alpha x+O(x^2)$ (Proposition~\ref{prop:tasks-aggregate}(a)), which gives $\Delta R_\sigma=x\sigma^*$. Exactly, $\Pi$ is proportional to $\ell^{e}$ by \eqref{eq:pb-closed} and $R^*_X=R_\infty+\Pi$ with $R_\infty$ free of $\lambda$ (Proposition~\ref{prop:saturation}(a)), so a factor $f$ on $\lambda$ changes $R^*_X$ by $\Pi(f^{e}-1)$: $f=1+x$ gives $\Pi[(1+x)^{e}-1]$, and $f=(1-x)^{-\alpha}$ gives $\Pi[(1-x)^{-\alpha e}-1]$. The function $\phi(x):=(1-x)^{-\alpha e}-1$ is convex with $\phi(0)=0$ and $\phi'(0)=\alpha e$, so $\phi(x)\ge\alpha e x$ and $\Pi\phi(x)\ge\Pi\alpha e x=(1-\alpha)(\sigma^*/\alpha)\alpha x/(1-\alpha)=x\sigma^*$, the inequality of Proposition~\ref{prop:tasks-aggregate}(b).
\emph{Overhead and engineering.} Both lines are part of $a$, which enters \eqref{eq:Rstar} only through $m_0=\gamma-a$ with coefficient one (Proposition~\ref{prop:envelope}); a cut $xa$ raises $R^*_X$ and $R_\infty$ by $xa$, exactly, since $R^*_X$ is affine in $m_0$.
\emph{Cost of revenue.} A cut $xc_h$ in human service cost raises $\gamma$ by $(1-\varphi)xc_h$ (Lemma~\ref{lem:prices}; Corollary~\ref{cor:priorities}(iv)) and so $R^*_X$ and $R_\infty$ by the same amount. When every rival makes the saving, the price falls at the rate $\bar\varphi^{\rm ind}$ and the margin is computed at the lower price, so $1-\varphi=1-\bar\varphi^{\rm ind}C_0$ to first order (Proposition~\ref{prop:tasks-aggregate}(a)); the volume response changes dollar profit (Corollary~\ref{cor:ext-game-engine}(a)) but not the margin share, because the other costs are per unit. Part (d) compares $x(1-\varphi)c_h$ with $xa_{\rm GA}$, or with $x(1-\varphi)a_{\rm GA}$ if the overhead cut is replicated too.
\emph{Retention.} $G_i$ enters $R^*_X$ only through $(X-k)G_i$ (Remark~\ref{rem:assumptions}(b) and \eqref{eq:Rstar}), so $\Delta R_G=y\,dG$ exactly; it also enters $R_\infty$ with the same coefficient.
\emph{Capital.} $\partial R^*_X/\partial k=-(g^*+\delta_K)$ with replacement at $\delta_K$ (Proposition~\ref{prop:envelope} and Remark~\ref{rem:assumptions}(c)), so $\Delta R_k=da-(g^*+\delta_K)dk$ to first order (Corollary~\ref{cor:capital}).
\emph{Durability.} $R_\infty=m_0+(X-k)G_i$ contains $m_0$ and $G_i$ and not $\lambda$, and $\Pi\to0$ as $s\to\bar s$ (Proposition~\ref{prop:saturation}(a), (b)).

(b) $\Delta R_\sigma-\Delta R_a=x[(\alpha\ell y)^{e}-a]$ has the sign of $\alpha\ell y-a^{1-\alpha}$, which is positive if and only if $y>a^{1-\alpha}/(\alpha\ell)$, that is, $X_{\mathsf o}>X^\ddagger_0$. Under Assumption~\ref{ass:playbook-calib}, $\alpha\ell=\sigma^{1-\alpha}/(\hat X-k)$ by Lemma~\ref{lem:playbook-calib}(a), so the condition reads $(X_{\mathsf o}-k)/(\hat X-k)>(a/\sigma)^{1-\alpha}$. For productivity, $\Delta R_\lambda>\Delta R_a$ if and only if $(\alpha\ell y)^{e}>\alpha a$, that is, $y>(\alpha a)^{1-\alpha}/(\alpha\ell)$. For the limit, \eqref{eq:reg-Xdag} with $da=\delta_\sigma a$ reads $X^\ddagger-k=\ell^{-1}[\delta_\sigma a/(C_\alpha\phi(\delta_\sigma))]^{1-\alpha}$; since $\phi(\delta_\sigma)/\delta_\sigma\to\alpha e$ and $C_\alpha\alpha e=\alpha^{e}$, the bracket tends to $a\alpha^{-e}$ and $X^\ddagger-k\to a^{1-\alpha}/(\alpha\ell)$. Finally, $\Delta R_\sigma$ and $\Delta R_\lambda$ are strictly increasing in $y$ by \eqref{eq:pb-closed}, $\Delta R_G$ is linear in $y$, and $\Delta R_a$, $\Delta R_h$ do not depend on $X$; so if $X^{\rm scr}<X^\ddagger_0<X_{\mathsf o}$ the overhead line ranks first at $X^{\rm scr}$ and the growth engine at $X_{\mathsf o}$.

(c) With $\ell=\lambda h(s)$, $\sigma^*=(\alpha\lambda h y)^{e}$ is proportional to $h^{e}$ at fixed $\lambda$ and $y$, and so are $\Delta R_\sigma$ and $\Delta R_\lambda$. At headroom $h'$ the growth engine outranks a line of size $a$ if and only if $\sigma^*(h'/h)^{e}>a$, that is, $h'>h(a/\sigma^*)^{1-\alpha}=h^\ddagger$. As $s\to\bar s$, $h\to0$ and both growth-engine values vanish, while the others do not depend on $h$ (part (a)).

(d) Compare $x\,\mathrm{dur}\,c_h$ and $x(1-\varphi)c_h$ with $xa_{\rm GA}$. Under linear demand $\bar\varphi^{\rm ind}=n_c/(n_c+1)$ with $n_c=1/\mathrm{HHI}$, so $\mathrm{dur}=1-(1-\mathrm{HHI})/(1+\mathrm{HHI})=2\mathrm{HHI}/(1+\mathrm{HHI})$, which tends to zero with $\mathrm{HHI}$; with constant curvature $\mathsf e$, $\mathrm{dur}=(2-\mathsf e)/(1/\mathrm{HHI}+1-\mathsf e)$, which is negative for $\mathsf e>2$ when $1/\mathrm{HHI}+1-\mathsf e>0$ (Corollary~\ref{cor:ext-game-engine}(a)). Under free entry without a scarce input the long-run markup is zero and the long-run score changes only through the return on capital (Proposition~\ref{prop:entry}(b), Corollary~\ref{cor:priorities}(iv)); with $k_u=0$ it does not change.

(e) By part (a), $m_0$ and $G_i$ enter $R^*_X$ only through $m_0+(X-k)G_i$, so the plan changes $R^*_{X_{\mathsf o}}$ by $-D_m+y\Delta G$ exactly, positive if and only if $D_m/\Delta G<y$ (Corollary~\ref{cor:prod-dilution}(i)); with $D_m=y$ and $\Delta G=1$ point the change is zero, so a point of retention is worth $y$ points of margin.

(f) The score part is Corollary~\ref{cor:capital} with replacement (Remark~\ref{rem:assumptions}(c)); the value part is Proposition~\ref{prop:tasks-capital}(a) at the owner's rate, per dollar of revenue; the sponsor part is Proposition~\ref{prop:reg-value}(d). The score hurdle $g^*+\delta_K$ increases in $g^*$.

(g) Proposition~\ref{prop:acquired-score}(a): $R^{\rm A}_X=R^{\rm org}_X+(X-k^{\rm acq})g^{\rm acq}$. \qed

\paragraph{Proof of Proposition~\ref{prop:playbook-gates}.}
Each item restates the cited result for the lever it gates; we record why the lever is worth nothing when the gate fails. (a), (b) An activity outside $\mathcal J^{\rm per}$ or before its feasibility date has no AI mode in $\mathcal M_j$, so its automation advantage is $-\infty$ (Definition~\ref{def:ext-effcost}) and every lever that works through its automation has no value. (c) Below the floor mode A is infeasible (Proposition~\ref{prop:tasks-review}(d)); below $\iota^*_p$ no activity of an all-or-nothing process is automated (Proposition~\ref{prop:tasks-readiness}(b)), so readiness investment that stops below $\iota^*_p$ buys nothing (part (d) of that proposition); a residency share above $\bar\omega^{\rm res}_p$ returns the process to its best non-automated mode (Proposition~\ref{prop:labor-residency}(b)). (d) If $e_j(1-d^R_j)>\bar\varepsilon_j$ no review share meets the target, so A is infeasible however cheap inference becomes (Propositions~\ref{prop:tasks-review}(d), \ref{prop:tasks-halluc}(b)); with automation bias the required catch rate is higher (Proposition~\ref{prop:tasks-bias}(d)). (e) The value of automation is $v_j[\underline c_j-c_j(A)]^+$ (Proposition~\ref{prop:tasks-modes}(c)), zero below the volume thresholds of Propositions~\ref{prop:tasks-modes}(e), \ref{prop:tasks-window}(c) and \ref{prop:tasks-buildbuy}(a) and at any inference price if the run-cost floor of Proposition~\ref{prop:tasks-fullcost}(b) fails; each threshold rises with the owner's rate (Proposition~\ref{prop:reg-speed}(a)). (f) The firm transitions if and only if the NPV is positive (Proposition~\ref{prop:tasks-transition}(b)), lowered for a financial process by the terms of Corollary~\ref{cor:acct-npv}; a pilot outside the band of \eqref{eq:tasks-pilotband} is worth less than scaling at once or doing nothing (Proposition~\ref{prop:tasks-pilot}(c)); past $r_DT_{\rm emb}=1$ a longer window lowers the expected value secured at embedding (Proposition~\ref{prop:ext-shocks}(b)). (g) Corollary~\ref{cor:labor-pace}, Propositions~\ref{prop:reg-dip}(b), \ref{prop:reg-speed}(b) and \ref{prop:labor-lead}(a): processes that cannot be completed in time lose the value the deadline or budget assigns them. (h) Corollary~\ref{cor:ext-game-engine}(b), Proposition~\ref{prop:ext-should}(b), (c) and Proposition~\ref{prop:acq-owner}(c). \qed

\paragraph{Proof of Corollary~\ref{cor:playbook-nopay}.}
(i) Proposition~\ref{prop:playbook-rank}(b), (c). (ii) Proposition~\ref{prop:playbook-rank}(d) and Corollary~\ref{cor:ext-game-engine}(b). (iii) Proposition~\ref{prop:playbook-rank}(f). (iv) Proposition~\ref{prop:playbook-gates}(c), (d). (v) Proposition~\ref{prop:playbook-gates}(e) and Proposition~\ref{prop:tasks-window}(a), (b). (vi) Proposition~\ref{prop:playbook-rank}(e) and Proposition~\ref{prop:prod-portfolio}(c). (vii) In \eqref{eq:acq-replace}, $C^{\rm cu}=X^{\mathsf o}S(1-e^{-\delta^GT^{\rm cu}})e^{-rt^{\rm rpl}}$ is proportional to $X^{\mathsf o}$ and increasing in $\delta^G$, so the condition $\mathcal G^{\rm spd}>\Delta\mathcal C$ fails for large enough $X^{\mathsf o}$ or $\delta^G$ (Proposition~\ref{prop:acq-replace}(a), (b)). (viii) Proposition~\ref{prop:playbook-rank}(g) and Proposition~\ref{prop:acq-limits}(a). (ix) Proposition~\ref{prop:playbook-gates}(f). \qed

\section{Playbook tables}\label{app:playbook-tables}

Table~\ref{tab:playbook-grid} gives the lever map for every archetype and owner, Table~\ref{tab:playbook-processes} the step decisions by process for archetype A2, and Table~\ref{tab:playbook-steps} every step of A2 at readiness 0.4. All three are generated by the archetype runner, \texttt{run\_archetypes.py} in \texttt{code/playbook}; the full records, for every archetype, readiness case, step and activity unit, are in \texttt{data/playbook} with their data dictionary. In Table~\ref{tab:playbook-grid} the public, dual-class and family weights are multiples of $\hat X-k$, so their rankings against $X^\ddagger_0=k+(\hat X-k)(a/\sigma)^{1-\alpha}$ do not depend on the level of $\hat X$; the sponsor, VC and permanent-acquirer weights are computed from rates, horizons, margins and exit prices, so their rankings are conditional on the level $\hat X-k=\mathsf r_\sigma/\alpha$ that the calibration implies (Lemma~\ref{lem:playbook-calib}(c)), and a lower level lowers $X^\ddagger_0$ and turns more of them toward the growth engine. % XBOUNDS

\IfFileExists{part2/tables/playbook_grid.tex}{\begin{table}[p]
\centering
\caption{Lever values by archetype and owner (points of the maximal score; the growth engine, overhead, engineering and cost of revenue per 10\% of the line, retention per point of zero-spend growth) and the engine's program value at readiness 0.4 and 0.9}\label{tab:playbook-grid}
\scriptsize
\setlength{\tabcolsep}{3pt}
\begin{tabular}{@{}llrrrrrrrrlr@{}}
\toprule
 & Owner & $X_{\mathsf o}$ & $\sigma^*$ & Growth & G\&A & R\&D & CoR & Retention & $X^\ddagger_0$ & First lever & Engine value (\$M) \\
\midrule
A1 & Public (preset pressure) & 3.4 & 24.6 & 2.5 & 2.4 & 3.2 & 1.0 & 3.3 & 3.3 & engineering (R\&D) & 4269 / 4310 \\
A1 & Dual class & 3.8 & 29.4 & 2.9 & 2.4 & 3.2 & 1.0 & 3.7 & 3.3 & engineering (R\&D) & 4270 / 4310 \\
A1 & Family & 4.4 & 35.9 & 3.6 & 2.4 & 3.2 & 1.0 & 4.3 & 3.3 & growth engine & 5671 / 5650 \\
A1 & PE, entry of a 6-year hold & 0.4 & 0.8 & 0.1 & 2.4 & 3.2 & 1.0 & 0.3 & 3.3 & engineering (R\&D) & 0 / 0 \\
A1 & VC, 7-year horizon & 1.7 & 9.2 & 0.9 & 2.4 & 3.2 & 1.0 & 1.7 & 3.3 & engineering (R\&D) & 1541 / 1578 \\
A1 & Permanent acquirer & 0.2 & 0.2 & 0.0 & 2.4 & 3.2 & 1.0 & 0.1 & 3.3 & engineering (R\&D) & 0 / 0 \\
\addlinespace
A2 & Public (preset pressure) & 3.4 & 1.4 & 0.1 & 1.0 & 2.0 & 0.7 & 3.3 & 13.1 & engineering (R\&D) & 5576 / 5647 \\
A2 & Dual class & 3.8 & 1.7 & 0.2 & 1.0 & 2.0 & 0.7 & 3.7 & 13.1 & engineering (R\&D) & 5578 / 5648 \\
A2 & Family & 4.4 & 2.1 & 0.2 & 1.0 & 2.0 & 0.7 & 4.3 & 13.1 & engineering (R\&D) & 7270 / 7303 \\
A2 & PE, entry of a 6-year hold & 3.3 & 1.3 & 0.1 & 1.0 & 2.0 & 0.7 & 3.2 & 13.1 & engineering (R\&D) & 1427 / 1484 \\
A2 & VC, 7-year horizon & 5.5 & 2.9 & 0.3 & 1.0 & 2.0 & 0.7 & 5.4 & 13.1 & engineering (R\&D) & 2481 / 2537 \\
A2 & Permanent acquirer & 5.4 & 2.8 & 0.3 & 1.0 & 2.0 & 0.7 & 5.3 & 13.1 & engineering (R\&D) & 983 / 1027 \\
\addlinespace
A3 & Public (preset pressure) & 3.5 & 0.0 & 0.0 & 1.0 & 1.5 & 0.7 & 3.3 & $\infty$ & engineering (R\&D) & 1449 / 1479 \\
A3 & Dual class & 3.9 & 0.0 & 0.0 & 1.0 & 1.5 & 0.7 & 3.7 & $\infty$ & engineering (R\&D) & 1452 / 1480 \\
A3 & Family & 22.4 & 0.0 & 0.0 & 1.0 & 1.5 & 0.7 & 22.3 & $\infty$ & engineering (R\&D) & 1897 / 1934 \\
A3 & PE, entry of a 6-year hold & 2.3 & 0.0 & 0.0 & 1.0 & 1.5 & 0.7 & 2.1 & $\infty$ & engineering (R\&D) & 589 / 632 \\
A3 & VC, 7-year horizon & 2.8 & 0.0 & 0.0 & 1.0 & 1.5 & 0.7 & 2.6 & $\infty$ & engineering (R\&D) & 635 / 661 \\
A3 & Permanent acquirer & 1.3 & 0.0 & 0.0 & 1.0 & 1.5 & 0.7 & 1.2 & $\infty$ & engineering (R\&D) & 277 / 302 \\
\addlinespace
A4 & Public (preset pressure) & 3.4 & 17.9 & 1.8 & 1.9 & 2.6 & 0.9 & 3.3 & 3.6 & engineering (R\&D) & 1533 / 1543 \\
A4 & Dual class & 3.8 & 21.3 & 2.1 & 1.9 & 2.6 & 0.9 & 3.7 & 3.6 & engineering (R\&D) & 1534 / 1544 \\
A4 & Family & 4.4 & 26.1 & 2.6 & 1.9 & 2.6 & 0.9 & 4.3 & 3.6 & engineering (R\&D) & 1881 / 1991 \\
A4 & PE, entry of a 6-year hold & 1.1 & 3.1 & 0.3 & 1.9 & 2.6 & 0.9 & 1.0 & 3.6 & engineering (R\&D) & 0 / 0 \\
A4 & VC, 7-year horizon & 1.6 & 5.6 & 0.6 & 1.9 & 2.6 & 0.9 & 1.5 & 3.6 & engineering (R\&D) & 668 / 683 \\
A4 & Permanent acquirer & 0.6 & 1.2 & 0.1 & 1.9 & 2.6 & 0.9 & 0.5 & 3.6 & engineering (R\&D) & 0 / 0 \\
\addlinespace
A5 & Public (preset pressure) & 4.5 & 6.2 & 0.6 & 2.0 & 2.5 & 0.9 & 3.3 & 8.5 & engineering (R\&D) & 283 / 306 \\
A5 & Dual class & 4.9 & 7.4 & 0.7 & 2.0 & 2.5 & 0.9 & 3.7 & 8.5 & engineering (R\&D) & 284 / 307 \\
A5 & Family & 5.9 & 10.4 & 1.0 & 2.0 & 2.5 & 0.9 & 4.7 & 8.5 & engineering (R\&D) & 360 / 373 \\
A5 & PE, entry of a 6-year hold & -0.1 & 0.0 & 0.0 & 2.0 & 2.5 & 0.9 & -1.3 & 8.5 & engineering (R\&D) & 0 / 0 \\
A5 & VC, 7-year horizon & 0.9 & 0.0 & 0.0 & 2.0 & 2.5 & 0.9 & -0.3 & 8.5 & engineering (R\&D) & 74 / 84 \\
A5 & Permanent acquirer & 0.1 & 0.0 & 0.0 & 2.0 & 2.5 & 0.9 & -1.1 & 8.5 & engineering (R\&D) & 0 / 0 \\
\addlinespace
A6 & Public (preset pressure) & 3.3 & 0.0 & 0.0 & 1.5 & 2.1 & 1.0 & 3.3 & $\infty$ & engineering (R\&D) & 3744 / 3774 \\
A6 & Dual class & 3.8 & 0.0 & 0.0 & 1.5 & 2.1 & 1.0 & 3.7 & $\infty$ & engineering (R\&D) & 3745 / 3774 \\
A6 & Family & 4.3 & 0.0 & 0.0 & 1.5 & 2.1 & 1.0 & 4.3 & $\infty$ & engineering (R\&D) & 4904 / 4930 \\
A6 & PE, entry of a 6-year hold & 11.3 & 0.0 & 0.0 & 1.5 & 2.1 & 1.0 & 11.2 & $\infty$ & engineering (R\&D) & 564 / 584 \\
A6 & VC, 7-year horizon & 17.6 & 0.0 & 0.0 & 1.5 & 2.1 & 1.0 & 17.5 & $\infty$ & engineering (R\&D) & 1106 / 1123 \\
A6 & Permanent acquirer & \multicolumn{10}{l}{retention growth G\_i at or above the unlevered hurdle} \\
\addlinespace
\bottomrule
\end{tabular}

\smallskip
\begin{minipage}{0.98\linewidth}\scriptsize
Notes: $X_{\mathsf o}$ at the archetype's state from Proposition~\ref{prop:reg-weight} and Proposition~\ref{prop:acq-pc}(d), evaluated with the engine's formulas and regime presets: public, the revealed weight $\hat X$; dual class, $X^{\rm my}$ at the dual-class pressure from the going-concern weight $k+(1+\mathfrak w)(\hat X-k)$; family, that weight scaled by $v_0(r_{\rm fam})/v_0(r_w)$; sponsor and VC, eq.~\eqref{eq:reg-Xexit} at entry with $\hat m=m_0-kG_i$ and, for the sponsor, the EBITDA rule at the marginal EBITDA margin $m_0$; permanent acquirer, $X^G=\hat m/(\bar{\mathfrak i}_{\rm U}-G_i)$, defined only for $G_i<\bar{\mathfrak i}_{\rm U}$; $\sigma^*$ is the owner's optimal growth spend (\% of revenue). Growth spend on Part~I's central, maintenance-netted basis (Assumption~\ref{ass:playbook-calib}). $X^\ddagger_0$ is the weight above which the growth engine outranks overhead (Proposition~\ref{prop:playbook-rank}(b)), infinite where S\&M is below the maintenance base and there is no growth spend. Engine value: the program value of the engine's own run for that regime (its presets) at readiness 0.4 / 0.9; its public, dual-class, sponsor and acquirer weights equal the $X_{\mathsf o}$ shown (\texttt{notes/playbook\_engine\_spec.md}, items W1, W14, W15, W17), while its VC weight uses its own exit rule; its sponsor and acquirer runs select no process for A1, A4 and A5; the engine has no permanent-acquirer platform, so that row values the firm's own transformation at the acquirer's hurdle. Generated by \texttt{code/playbook/run\_archetypes.py}.
\end{minipage}
\end{table}
}{}
\IfFileExists{part2/tables/playbook_processes.tex}{\begin{table}[p]
\centering
\caption{Step decisions by process for archetype A2: modes of the steps in the engine's best process plan starting now (counts by mode) at readiness 0.4 and 0.9, the most common binding constraint on automation, ICFR steps, the wave and the tiers}\label{tab:playbook-processes}
\scriptsize
\setlength{\tabcolsep}{3pt}
\begin{tabular}{@{}p{4.0cm}rllp{4.0cm}rrl@{}}
\toprule
Process & Steps & Modes, 0.4 & Modes, 0.9 & Binding constraint (most steps) & ICFR & Wave & Tiers\\
\midrule
P01 Order to cash & 6 & H1O2U3 & H1O2U1N2 & feasibility: capability horizon & 6 & -- & AC \\
P02 Procure to pay & 6 & O3U3 & O2U2N2 & readiness & 6 & -- & ABC \\
P03 Record to report & 7 & H2O5 & H2O5 & feasibility: capability horizon & 7 & -- & AB \\
P04 Month-end close & 5 & H5 & H5 & feasibility: physical & 5 & -- & AB \\
P05 Hire to retire & 7 & H2O5 & H1O5U1 & feasibility: capability horizon & 3 & -- & ABC \\
P06 Payroll & 5 & O5 & O5 & feasibility: capability horizon & 5 & -- & ABC \\
P07 Ticket to resolution & 6 & H1O3U2 & H1O2U2N1 & readiness & 1 & -- & ABC \\
P08 Lead to close & 7 & H1O1U5 & H1U4N2 & readiness & 1 & 1 & AB \\
P09 Quote to contract & 5 & O1U4 & U3N2 & readiness & 1 & -- & AB \\
P10 Idea to release & 7 & O4U3 & O1U4N2 & feasibility: capability horizon & 2 & -- & A \\
P11 Incident to resolution & 5 & H1O3U1 & H1O1U1N2 & feasibility: capability horizon & 1 & -- & AB \\
P12 Campaign to lead & 5 & H2U3 & H2U3 & feasibility: capability horizon & 0 & -- & ABC \\
P13 Onboard to renew & 4 & H1O1U2 & H1O1U2 & feasibility: physical & 2 & -- & ABC \\
P14 Implementation project & 6 & O2U4 & U4N2 & feasibility: capability horizon & 2 & -- & ABC \\
P15 Contract review and legal advice & 5 & O2U3 & O1U3N1 & permissibility & 1 & -- & AC \\
P16 Control test to remediation & 4 & H2O2 & H2O1N1 & permissibility & 4 & -- & AC \\
P17 IT service request & 5 & O5 & O2U2N1 & feasibility: capability horizon & 2 & -- & AC \\
P18 Plan to budget & 6 & H3O3 & H3O3 & feasibility: capability horizon & 0 & -- & AB \\
P19 Expense to reimburse & 4 & H1O3 & H1O3 & feasibility: physical & 4 & -- & A \\
P20 Order fulfilment & 6 & H2O4 & H2O3N1 & readiness & 3 & -- & AB \\
P21 Plan to produce & 6 & H2O3U1 & H2O3N1 & feasibility: physical & 2 & -- & AB \\
P22 Facilities work order & 3 & H1O2 & H1O1N1 & feasibility: physical & 0 & -- & A \\
P23 Strategy to plan & 5 & H5 & H5 & feasibility: capability horizon & 0 & -- & A \\
P24 Tax compliance & 5 & H5 & H5 & readiness & 5 & -- & AB \\
\bottomrule
\end{tabular}

\smallskip
\begin{minipage}{0.98\linewidth}\scriptsize
Notes: H in house, O outsourced, U augmented, A automated, N eliminated; a step keeps its baseline mode when no column of its process has positive value. Generated by \texttt{code/playbook/run\_archetypes.py}; rows in \texttt{data/playbook/playbook\_steps.csv}.
\end{minipage}
\end{table}
}{}
\IfFileExists{part2/tables/playbook_steps_full.tex}{\begingroup\scriptsize
\setlength{\tabcolsep}{3pt}
\begin{longtable}{@{}lp{4.4cm}llp{4.2cm}rll@{}}
\caption{Archetype A2 at step level, readiness 0.4: activity and step decisions, the mode at readiness 0.9, the binding constraint on automation, the detection gate where reliability binds, ICFR scope and the confidence tier}\label{tab:playbook-steps}\\
\toprule
Step & Name & Now (act./step) & 0.9 & Binding constraint & $d^*$ & ICFR & Tier\\
\midrule
\endfirsthead
\toprule
Step & Name & Now (act./step) & 0.9 & Binding constraint & $d^*$ & ICFR & Tier\\
\midrule
\endhead
\bottomrule
\endfoot
P01.1 & Verify customer and order data & U/U & U & permissibility &  & yes & C \\
P01.2 & Prepare invoice & U/U & N & feasibility: capability horizon &  & yes & C \\
P01.3 & Record receivable & U/U & N & feasibility: capability horizon &  & yes & C \\
P01.4 & Collect payment & H/H & H & feasibility: physical &  & yes & A \\
P01.5 & Apply cash and reconcile & O/O & O & feasibility: physical &  & yes & A \\
P01.6 & Resolve billing disputes & O/O & O & readiness &  & yes & C \\
P02.1 & Identify need and gather supplier information & O/O & O & readiness &  & yes & C \\
P02.2 & Negotiate terms & U/U & U & readiness &  & yes & C \\
P02.3 & Issue purchase order & U/U & U & readiness &  & yes & B \\
P02.4 & Receive and inspect goods or services & \multicolumn{6}{l}{not in the activity map of this firm} \\
P02.5 & Match invoice to order and receipt & O/O & O & readiness &  & yes & B \\
P02.6 & Execute payment & O/O & N & feasibility: physical &  & yes & A \\
P02.7 & Record payable & U/U & N & readiness &  & yes & C \\
P03.1 & Compile ledger and subledger data & O/O & O & feasibility: capability horizon &  & yes & A \\
P03.2 & Post journal entries & O/O & O & feasibility: capability horizon &  & yes & A \\
P03.3 & Reconcile accounts & O/O & O & feasibility: physical &  & yes & A \\
P03.4 & Analyze variances & O/O & O & feasibility: capability horizon &  & yes & A \\
P03.5 & Prepare financial statements & O/O & O & feasibility: capability horizon &  & yes & A \\
P03.6 & Review for accuracy and compliance & H/H & H & permissibility &  & yes & A \\
P03.7 & Explain results to management & H/H & H & readiness &  & yes & B \\
P04.1 & Calculate accruals and adjustments & O/H & O & readiness &  & yes & B \\
P04.2 & Reconcile balance sheet accounts & O/H & O & feasibility: physical &  & yes & A \\
P04.3 & Consolidate entities & O/H & O & feasibility: capability horizon &  & yes & A \\
P04.4 & Review close checklist and controls & H/H & H & feasibility: physical &  & yes & A \\
P04.5 & Prepare close package & O/H & O & feasibility: capability horizon &  & yes & A \\
P05.1 & Plan staffing needs & H/H & H & readiness &  &  & C \\
P05.2 & Recruit candidates & O/O & O & readiness &  &  & B \\
P05.3 & Interview candidates & O/O & O & feasibility: capability horizon &  &  & C \\
P05.4 & Evaluate and select candidates & H/H & U & feasibility: capability horizon &  &  & A \\
P05.5 & Onboard and train & O/O & O & feasibility: capability horizon &  & yes & C \\
P05.6 & Administer HR records and benefits & O/O & O & feasibility: capability horizon &  & yes & C \\
P05.7 & Review performance & \multicolumn{6}{l}{not in the activity map of this firm} \\
P05.8 & Process separation & O/O & O & readiness &  & yes & B \\
P06.1 & Compile time and attendance & O/O & O & feasibility: capability horizon &  & yes & C \\
P06.2 & Calculate pay and deductions & O/O & O & readiness &  & yes & C \\
P06.3 & Execute payment & O/O & O & feasibility: physical &  & yes & A \\
P06.4 & Record payroll & O/O & O & feasibility: capability horizon &  & yes & A \\
P06.5 & File payroll taxes & O/O & O & readiness &  & yes & B \\
P07.1 & Receive and triage inquiry & O/O & O & readiness &  &  & C \\
P07.2 & Gather customer information & O/O & O & readiness &  &  & C \\
P07.3 & Diagnose problem & H/H & H & feasibility: physical &  &  & A \\
P07.4 & Explain solution & U/U & U & readiness &  &  & C \\
P07.5 & Escalate and coordinate fix & U/U & A & readiness &  & yes & B \\
P07.6 & Record resolution & O/O & N & feasibility: capability horizon &  &  & A \\
P08.1 & Prospect and build relationships & U/U & U & readiness &  &  & B \\
P08.2 & Qualify needs & U/U & N & feasibility: capability horizon &  &  & A \\
P08.3 & Present and demonstrate & U/U & U & readiness &  &  & B \\
P08.4 & Price the offer & U/U & U & readiness &  &  & B \\
P08.5 & Negotiate & U/U & U & readiness &  &  & B \\
P08.6 & Close sale & H/H & H & feasibility: physical &  & yes & A \\
P08.7 & Record in CRM & O/O & N & feasibility: capability horizon &  &  & A \\
P09.1 & Capture requirements & U/U & N & feasibility: capability horizon &  &  & A \\
P09.2 & Prepare quote & U/U & U & readiness &  &  & B \\
P09.3 & Draft contract & O/O & N & feasibility: capability horizon &  &  & A \\
P09.4 & Review terms for compliance & \multicolumn{6}{l}{not in the activity map of this firm} \\
P09.5 & Negotiate terms & U/U & U & readiness &  &  & B \\
P09.6 & Approve and sign & U/U & U & readiness &  & yes & B \\
P10.1 & Gather requirements & U/U & N & feasibility: capability horizon &  &  & A \\
P10.2 & Design solution & U/U & U & feasibility: capability horizon &  &  & A \\
P10.3 & Write code & O/O & U & feasibility: capability horizon &  &  & A \\
P10.4 & Test & U/U & U & feasibility: capability horizon &  & yes & A \\
P10.5 & Document & O/O & N & feasibility: capability horizon &  &  & A \\
P10.6 & Deploy & O/O & O & feasibility: physical &  & yes & A \\
P10.7 & Monitor in production & O/O & U & feasibility: capability horizon &  &  & A \\
P11.1 & Detect through monitoring & O/O & U & feasibility: capability horizon &  &  & A \\
P11.2 & Notify responders & U/U & A & readiness &  &  & B \\
P11.3 & Diagnose & H/H & H & feasibility: physical &  &  & A \\
P11.4 & Restore service & O/O & O & feasibility: physical &  & yes & A \\
P11.5 & Write post-incident report & O/O & N & feasibility: capability horizon &  &  & A \\
P12.1 & Analyze market & U/U & U & feasibility: capability horizon &  &  & C \\
P12.2 & Plan campaign & H/H & H & feasibility: capability horizon &  &  & A \\
P12.3 & Create content & U/U & U & feasibility: capability horizon &  &  & A \\
P12.4 & Run promotion & H/H & A & readiness &  &  & B \\
P12.5 & Measure response & U/U & U & feasibility: capability horizon &  &  & C \\
P13.1 & Train customer users & O/O & O & feasibility: physical &  &  & A \\
P13.2 & Advise on adoption & U/U & U & readiness &  &  & B \\
P13.3 & Resolve account issues & U/U & U & readiness &  & yes & C \\
P13.4 & Renew and expand & H/H & H & feasibility: physical &  & yes & A \\
P14.1 & Confer with client on scope & U/U & N & feasibility: capability horizon &  &  & C \\
P14.2 & Plan project & U/U & U & feasibility: capability horizon &  & yes & A \\
P14.3 & Configure systems & O/O & U & feasibility: capability horizon &  & yes & A \\
P14.4 & Train client staff & U/U & U & feasibility: capability horizon &  &  & A \\
P14.5 & Document solution & O/O & N & feasibility: capability horizon &  &  & A \\
P14.6 & Coordinate with client and partners & U/U & U & readiness &  &  & B \\
P15.1 & Research law and precedent & O/O & O & permissibility &  &  & C \\
P15.2 & Review documents & U/U & U & readiness &  &  & C \\
P15.3 & Prepare legal documents & O/O & N & feasibility: capability horizon &  &  & A \\
P15.4 & Advise client & U/U & U & permissibility &  &  & C \\
P15.5 & Negotiate & U/U & U & readiness &  & yes & C \\
P16.1 & Plan test & \multicolumn{6}{l}{not in the activity map of this firm} \\
P16.2 & Examine evidence & H/H & H & permissibility &  & yes & C \\
P16.3 & Evaluate compliance & O/O & O & feasibility: capability horizon &  & yes & A \\
P16.4 & Document findings & O/O & N & feasibility: capability horizon &  & yes & A \\
P16.5 & Coordinate remediation & H/H & H & permissibility &  & yes & C \\
P17.1 & Receive request & O/O & U & feasibility: capability horizon &  &  & A \\
P17.2 & Provision access or system & O/O & O & feasibility: physical &  & yes & A \\
P17.3 & Apply security controls & O/O & U & feasibility: capability horizon &  & yes & A \\
P17.4 & Record ticket & O/O & N & feasibility: capability horizon &  &  & A \\
P17.5 & Instruct user & O/O & O & feasibility: capability horizon &  &  & C \\
P18.1 & Gather operating data & O/O & O & readiness &  &  & B \\
P18.2 & Analyze financial data & O/O & O & feasibility: capability horizon &  &  & A \\
P18.3 & Develop financial plan & H/H & H & feasibility: capability horizon &  &  & A \\
P18.4 & Prepare budget & O/O & O & feasibility: capability horizon &  &  & A \\
P18.5 & Manage budget & H/H & H & feasibility: capability horizon &  &  & A \\
P18.6 & Explain to management & H/H & H & readiness &  &  & B \\
P19.1 & Check receipts and policy & H/H & H & feasibility: physical &  & yes & A \\
P19.2 & Reconcile to card feed & O/O & O & feasibility: physical &  & yes & A \\
P19.3 & Reimburse & O/O & O & feasibility: physical &  & yes & A \\
P19.4 & Record expense & O/O & O & feasibility: capability horizon &  & yes & A \\
P20.1 & Schedule order & O/O & N & readiness &  &  & B \\
P20.2 & Check inventory & H/H & H & readiness &  & yes & B \\
P20.3 & Pick and move goods & O/O & O & feasibility: physical &  &  & A \\
P20.4 & Pack & O/O & O & feasibility: physical &  &  & A \\
P20.5 & Ship and document & O/O & O & readiness &  & yes & B \\
P20.6 & Replenish stock & H/H & H & readiness &  & yes & B \\
P21.1 & Schedule production & O/O & O & readiness &  &  & B \\
P21.2 & Set up equipment & O/O & O & feasibility: physical &  &  & A \\
P21.3 & Operate production equipment & H/H & H & feasibility: physical &  &  & A \\
P21.4 & Inspect output & H/H & H & feasibility: physical &  & yes & A \\
P21.5 & Maintain equipment & O/O & O & feasibility: physical &  &  & A \\
P21.6 & Record production & U/U & N & readiness &  & yes & B \\
P22.1 & Receive request & \multicolumn{6}{l}{not in the activity map of this firm} \\
P22.2 & Inspect site & H/H & H & feasibility: physical &  &  & A \\
P22.3 & Schedule work & \multicolumn{6}{l}{not in the activity map of this firm} \\
P22.4 & Repair or maintain & O/O & O & feasibility: physical &  &  & A \\
P22.5 & Record completion & O/O & N & feasibility: capability horizon &  &  & A \\
P23.1 & Analyze market and competitors & H/H & H & feasibility: capability horizon &  &  & A \\
P23.2 & Set goals & H/H & H & feasibility: capability horizon &  &  & A \\
P23.3 & Develop business plan & H/H & H & feasibility: capability horizon &  &  & A \\
P23.4 & Communicate strategy & H/H & H & feasibility: capability horizon &  &  & A \\
P23.5 & Direct execution & H/H & H & feasibility: capability horizon &  &  & A \\
P24.1 & Compile tax data & O/H & O & feasibility: capability horizon &  & yes & A \\
P24.2 & Calculate liability & O/H & O & readiness &  & yes & B \\
P24.3 & Prepare return & O/H & O & readiness &  & yes & B \\
P24.4 & Review & H/H & H & permissibility &  & yes & A \\
P24.5 & Pay and file & O/H & O & permissibility &  & yes & A \\
\end{longtable}
\endgroup
}{}
